\chapter{Generation Prompts}
\label{ch:prompts}

This chapter turns the document into code. Each prompt is written to be pasted into an agent
session together with this document. They assume the agent can read files, run commands and
create commits, and they assume nothing else.

\section{Rules of engagement}

Prepend this to every session. It is the contract that keeps generated work consistent with
the specification.

\begin{lstlisting}
You are implementing Xyra, an agentic IDE, from a written specification.

The specification is authoritative. If your instinct disagrees with it, follow the
specification and state the disagreement in your final message rather than silently
deviating.

Rules that apply to every session:
1. Implement exactly one issue identifier from the register. Do not start another.
2. Never add a dependency whose license is outside MIT, Apache-2.0, BSD, ISC, Zlib,
   Unlicense or CC0. If you need one that is not, stop and report it.
3. Write no code comments. Names carry the meaning. Rationale goes in the commit body.
4. Never use an em dash anywhere: not in code, strings, documents or commit messages.
5. Interface copy is sentence case. No capitalised words.
6. Every control you add must already exist in the interface chapter. If it does not,
   stop and report the gap rather than inventing a control.
7. Every long running operation takes a budget and a cancellation token.
8. Write tests, including at least one for the failure path. Run them. Report the result
   with the exact command and its exit code.
9. Leave the working tree building and the test suite green. Do not commit; report the
   diff and let the supervisor commit.
10. Never run an interactive command or a process that does not exit. Use the
    non interactive flag, or start it detached and read the log.

When you finish, report: the issue identifier, the files changed, the test command and
its result, and anything in the specification you found ambiguous.
\end{lstlisting}

\section{Bootstrap prompt}

Run once, on an empty repository.

\begin{lstlisting}
Read the attached blueprint. Create the repository skeleton described in the architecture
chapter: a workspace manifest, one directory per crate listed in the crate topology table,
a lint and format configuration, a license allow list check, and a continuous integration
definition that runs format, lint, test and the license check.

Every crate is created with a doc header stating its single responsibility as written in
the crate topology table, and with a placeholder test so the suite is green from the first
commit. Do not implement behaviour yet.

This is issue PLA-001. Follow the rules of engagement.
\end{lstlisting}

\section{Category prompts}

Each of the twelve categories has a lead prompt that establishes shared understanding before
individual issues are executed against it.

\subsection{Platform}
\begin{lstlisting}
You own the PLA category: the window shell, the element toolkit, the core service and
the transport between them.

Read the architecture chapter and the design system section of the interface chapter.
Implement the design tokens exactly as tabulated: spacing, radius, control heights, text
sizes, icon sizes. No arbitrary values anywhere in the toolkit.

Then implement the controls listed in PLA-004 through PLA-012, each with its full state
set: default, hover, active, focused, disabled, loading. Every control gets an accessible
name and keyboard reachability. Every control gets an interface test through the harness.

Work issue by issue in the order given in the register. One issue, one report.
\end{lstlisting}

\subsection{Text core}
\begin{lstlisting}
You own the TXT category: the rope buffer and everything that edits it.

Correctness is the entire job here. Implement transactions such that undo restores both
content and selection, coalescing is predictable, and concurrent transactions on one buffer
are ordered rather than merged optimistically.

Write property tests before implementation: random edit sequences must round trip through
undo and redo, and the rope must agree with a naive string model on every operation.

Start with TXT-001 and proceed in register order.
\end{lstlisting}

\subsection{Syntax}
\begin{lstlisting}
You own the SYN category: parsing and everything derived from a parse tree.

The parse tree is the substrate for the context engine, not only for colouring. Symbol
extraction in SYN-008 must produce, for every file, the ordered symbol list with full
signatures, type declarations, documentation comments and outbound call edges, in the
shape the context chapter calls a Structural Context Protocol record.

Design that output format first, write golden tests for it against fixture files in each
supported language, then implement.
\end{lstlisting}

\subsection{Language}
\begin{lstlisting}
You own the LNG category: language server integration.

Supervision is as important as the features. A server that crashes restarts with backoff
and the editor never stops working. A server that lacks a capability degrades that feature
visibly rather than failing silently.

Implement LNG-001 and LNG-002 fully before any feature issue, because every feature depends
on correct capability negotiation.
\end{lstlisting}

\subsection{Project and version control}
\begin{lstlisting}
You own the PRJ category: the worktree, watching and git.

Two invariants: the developer's edit always wins over any automated change, and no
operation may leave the repository in a state the developer did not ask for. Every
destructive operation is reversible or confirmed.

Implement staging by hunk in PRJ-007 with the same data model the diff view uses, so the
two surfaces cannot disagree.
\end{lstlisting}

\subsection{Terminal}
\begin{lstlisting}
You own the TRM category: terminal emulation and process supervision.

TRM-007 is the issue that matters most: enforce the non interactive policy on commands
invoked by agents. A command that would prompt is refused with the correct non interactive
form stated, a process that does not exit is started detached with its log tailed, and a
command producing no output for a bounded period is terminated.

Implement TRM-001 and TRM-002 first, then TRM-007 before the remaining features.
\end{lstlisting}

\subsection{Context engine}
\begin{lstlisting}
You own the CTX category: the structural graph, the compression protocol, the budget
solver and the retrieval pipeline.

Read the context chapter completely before writing code. Implement in this order:
structural extraction, then the SCP record at three levels, then token counting with a
real tokeniser, then the budget solver, then delta encoding, then the assembly pipeline.
Similarity search comes last and is always labelled as a fallback in its results.

Build the evaluation harness described in CTX-040 early, not at the end. Every subsequent
change to the engine must be justified by a measured improvement in first attempt success,
compression ratio or escalation rate against the naive baseline.
\end{lstlisting}

\subsection{Agent layer}
\begin{lstlisting}
You own the AGT category: agent sessions, the tool surface and skills.

Every tool has a typed contract, a capability gate and an audit record. No tool grants
ambient authority. Write access is a property of the role as tabulated in the agent
chapter, not of the model.

AGT-017, the inactivity watchdog and hard kill path, is not optional and is not last. A
hung agent process must never be able to block the interface or a mission.
\end{lstlisting}

\subsection{Mission}
\begin{lstlisting}
You own the MIS category: the supervisor that runs a project to completion unattended.

The state machine is specified in the agent chapter. Implement it exactly, including the
four step failure escalation: retry, vendor swap, split with bounded depth, quarantine and
continue.

Two hazards to handle explicitly. The working tree reset between attempts must spare the
supervisor's own state directory, otherwise the supervisor destroys its own memory. And
ticket splitting must be depth bounded, otherwise a persistently failing ticket splits
forever.

State is written atomically after every transition. Prove resumption with a test that kills
the process mid ticket and continues.
\end{lstlisting}

\subsection{Verification}
\begin{lstlisting}
You own the VER category: everything that proves a change works.

Verification never touches the developer's working tree. Snapshot the change set into a
throwaway worktree and execute there.

Implement VER-001 through VER-003 first, because the mission category depends on them.
Then targeted test selection, then visual verification, then the adversarial sweep.

Every verdict is structured, with severity ranked issues, because verdicts feed automatic
fix loops rather than human reading.
\end{lstlisting}

\subsection{Fleet}
\begin{lstlisting}
You own the FLT category: multi repository operations.

Registration comes from the developer's provider, never from a hand written file. Impact
analysis must separate the repository that defines a contract from those that consume it,
because that ordering is what makes a cross repository change safe.

Each repository is verified independently before any linked change is prepared.
\end{lstlisting}

\subsection{Extensions}
\begin{lstlisting}
You own the EXT category: the extension host.

Extensions are untrusted. No ambient filesystem access, no network unless granted, an
explicit capability list in the manifest, and a trap disables the extension for the
session without affecting the editor.

Implement EXT-002 and EXT-010 before any contribution point, because isolation is the
feature.
\end{lstlisting}

\section{Issue prompt template}

For each issue in the register, substitute and run.

\begin{lstlisting}
Implement issue {ID}: {title}.

Read these sections of the blueprint before writing code:
- the architecture chapter section for {crate}
- the interface chapter section for {surface}, if this issue touches the interface
- the user stories that reference this capability: {story ids}

Acceptance:
- The behaviour matches the specification exactly.
- Tests exist including a failure path test, and they pass.
- No control exists that is not specified in the interface chapter.
- No dependency was added outside the license allow list.
- The operation has a budget and a cancellation path if it can run long.
- Accessible name and keyboard reachability exist for every new control.

Dependencies already landed: {depends}.
Follow the rules of engagement. Report when done.
\end{lstlisting}

\section{Self hosting prompt}

The most efficient way to build Xyra is to use the existing Xyra mission supervisor to build
it. Once the register is loaded as a ticket graph, the loop in the agent chapter applies to
this project as it does to any other.

\begin{lstlisting}
Convert the issue register in the blueprint into a mission ticket graph. Each issue becomes
one ticket: the identifier is the ticket identifier, the issue text is the scope, and the
Depends column becomes the dependency list.

Then run the mission to completion with these budgets: attempts per ticket three, split
depth one, and a stop on five consecutive failures.

For every ticket, the brief is the issue prompt template filled in from the register, with
the rules of engagement prepended.

Verification for this project is the workspace test suite plus the format and lint checks
plus the license allow list check. A ticket counts as done only when all four pass in an
isolated snapshot.
\end{lstlisting}

\section{Review prompt}

Run against every diff before it is accepted, ideally on a different vendor from the one that
wrote it.

\begin{lstlisting}
You are reviewing a diff produced by another vendor against a written specification.

Check, in this order:
1. Correctness. Does it do what the issue says, including the failure path.
2. Specification fidelity. Does any control, message or behaviour deviate from the
   blueprint. Quote the deviation.
3. Boundaries. Does it call across a layer in the wrong direction, or reach into another
   category's crate.
4. Security. Untrusted input treated as instruction, secrets in context, missing capability
   gate, unbounded resource use.
5. Conventions. Comments present, em dashes present, capitalised interface copy, arbitrary
   spacing values instead of tokens.

Return a structured verdict: severity per finding, file and line, and a blocking decision.
Default to blocking when uncertain about correctness or security.
\end{lstlisting}

\section{Specification amendment prompt}

The document is the source of truth, which means it must be maintained as the code teaches
you things.

\begin{lstlisting}
While implementing {ID} you found that the blueprint is wrong, ambiguous or incomplete.

Do not silently deviate. Produce an amendment:
- The exact section and sentence that is wrong.
- What the code must do instead, and why the specification could not have anticipated it.
- Any other section that must change for the document to stay internally consistent.
- The stories and issues affected.

Emit the amendment as a patch to the blueprint source alongside the code change. The
document and the code land together or not at all.
\end{lstlisting}
